\documentclass[12pt]{amsart}


\usepackage{amssymb}
\usepackage{amsthm}
\usepackage{amsmath}
\usepackage{amsxtra}
\usepackage{mathrsfs}
\usepackage{enumerate}

\setlength{\hfuzz}{4pt}

\newenvironment{problab}[1]
{\noindent\textbf{Problem #1}.}
{\vskip 6pt}

\newenvironment{exolab}[1]
{\noindent\textbf{Exercise #1}.}
{\vskip 6pt}


\theoremstyle{remark}
\newtheorem*{soln}{Solution}

\newcommand{\C}{\mathbb C}
\newcommand{\D}{\mathbb D}
\newcommand{\F}{\mathbb F}
\newcommand{\HH}{\mathbb H}
\newcommand{\PP}{\mathbb P}
\newcommand{\Q}{\mathbb Q}
\newcommand{\R}{\mathbb R}
\newcommand{\bbS}{\mathbb S}
\newcommand{\Z}{\mathbb Z}

\newcommand{\eps}{\epsilon}

\usepackage{fullpage}

\DeclareMathOperator{\Aut}{Aut}
\DeclareMathOperator{\Hom}{Hom}
\DeclareMathOperator{\impart}{Im}
\DeclareMathOperator{\repart}{Re}

\usepackage{mathpazo}

\begin{document}

\title[MATH 351 Homework]{Belyi Maps and Dessins d'Enfants \\ Homework \#1}
\date{\today}

%\thanks{\emph{Date:} Monday, 11 February 2013}

\maketitle

\begin{exolab}{1.1}
  Let
  \begin{equation*}
    X = \{(x,y,z) \in \R^3 : x^2 + y^2 = z^2\}
  \end{equation*}
  be the cone in $\R^3$ equipped with the subspace topology. Show that $X$ cannot be given the structure of a topological surface.
\end{exolab}

\begin{exolab}{1.2}
  Fix $\omega_1, \omega_2 \in \C$ that are $\R$-linearly independent, and let $\Lambda = \Z \omega_1 \oplus \Z \omega_2$ be the lattice they span. Let $X$ be the quotient $\C/\Lambda$, a complex torus. Show that the collection of charts defined in lecture form a holomorphic atlas. (All that is left to check is that the transition functions are holomorphic.)
\end{exolab}

\begin{exolab}{1.3}
  \begin{enumerate}[(a)]
    \item 
      Let $z = a + bi \in \C$. Considering $\C$ as an $\R$-vector space, then the left-multiplication map
      \begin{align*}
        \lambda_z: \C &\to \C\\
        w &\mapsto zw
      \end{align*}
      is an $\R$-linear map. Determine the matrix of $\lambda_z$ with respect to the basis $\{1, i\}$.

    \item
      Let $h: \R^2 \to \R^2$ be a differentiable function given by component functions $h_1, h_2: \R^2 \to \R$, so $h(x,y) = (h_1(x,y), h_2(x,y))$. Recall that its derivative is given by the Jacobian matrix
      \begin{equation*}
        \begin{pmatrix}
          \frac{\partial h_1}{\partial x} & \frac{\partial h_1}{\partial y}\\[0.5em]
          \frac{\partial h_2}{\partial x} & \frac{\partial h_2}{\partial y}\\
        \end{pmatrix}
      \end{equation*}

      Let $U \subseteq \C$ be an open set and let $f: U \to \C$ be holomorphic. Then $f(z) = u(x,y) + iv(x,y)$ is differentiable considered as a function $\R^2 \to \R^2$, and its Jacobian matrix at $z_0$ must be the matrix of left multiplication by the complex number $f'(z_0)$.

      Recover the Cauchy-Riemann equations by equating the Jacobian matrix  and the matrix representation of $a+bi$ you found in the previous part.
  \end{enumerate}
\end{exolab}

\end{document}
